% !TEX program = lualatex
% Auto-generated from YAML (v3) — 04: 課題の見つけ方と解決プロセス

\documentclass[aspectratio=43,professionalfonts,handout,10pt]{beamer}
\usepackage{beamer_template}

\newcommand{\LectureNum}{04}
\newcommand{\LectureTitle}{課題の見つけ方と解決プロセス}
\newcommand{\CourseName}{情報リテラシー}
\newcommand{\TextPages}{03-05}

\begin{document}
% --- Slide 1: title ---

\begin{frame}{\CourseName\quad 第\LectureNum 回「\LectureTitle 」}
  {\small \TermName\quad \TeacherName\quad ／\quad 教科書 \TextPages}

  \vfill

  \begin{block}{今日の流れ}
    \begin{enumerate}
      \item 問題解決の5ステップ
      \item PDCAサイクル
      \item インターネットによる情報収集
      \item 発想・整理の手法（ブレーンストーミング・KJ法・マインドマップ）
    \end{enumerate}
  \end{block}
  \vfill

  \begin{block}{今日の目標}
    \begin{itemize}
      \item 問題解決の基本的なプロセスを説明できる。
    \end{itemize}
  \end{block}
  \vfill

  \begin{exampleblock}{キーワード}
    \small \textbf{問題（Problem）} ／ \textbf{課題（Task）} ／ \textbf{フィードバック} ／ \textbf{検索エンジン} ／ \textbf{ブレーンストーミング} ／ \textbf{KJ法} ／ \textbf{マインドマップ}
  \end{exampleblock}
\end{frame}

% --- Slide 2: 問題解決の5ステップ ---

\begin{frame}{問題解決の5ステップ（Theory 03）}
  \begin{block}{ポイント}
    \begin{itemize}
      \item ① 問題の発見：現状と理想のギャップを明確にする
      \item ② 情報の収集：問題に関連するデータ・事例・意見を集める
      \item ③ 情報の分析：収集した情報を整理・比較・評価する
      \item ④ 解決策の設計：分析結果をもとに具体的な解決策を考える
      \item ⑤ 解決策の実施：解決策を実行し，結果を評価してフィードバックする
    \end{itemize}
  \end{block}
\end{frame}

% --- Slide 3: PDCAサイクル ---

\begin{frame}{PDCAサイクル}
  \begin{block}{ポイント}
    \begin{itemize}
      \item Plan（計画）：目標を設定し，達成するための計画を立てる
      \item Do（実行）：計画に従って実際に行動する
      \item Check（評価）：実行結果を目標と比較して評価する
      \item Act（改善）：評価結果をもとに計画・行動を改善する
    \end{itemize}
  \end{block}

  \vfill

  \begin{exampleblock}{補足}
    \begin{itemize}
      \item PDCAは1回で終わらず，繰り返すことで継続的に改善していく仕組み
    \end{itemize}
  \end{exampleblock}

  \vfill

  \begin{alertblock}{演習}
    PDCAサイクルを学習に応用する例を挙げよ
  \end{alertblock}
\end{frame}

% --- Slide 4: インターネットによる情報収集 ---

\begin{frame}{インターネットによる情報収集（Theory 04）}
  \begin{block}{ポイント}
    \begin{itemize}
      \item AND検索：複数のキーワードをすべて含むページを検索（例：「情報 AND セキュリティ」）
      \item OR検索：いずれかのキーワードを含むページを検索（例：「情報 OR データ」）
      \item NOT検索：特定のキーワードを除いて検索（例：「情報 NOT 個人」）
      \item 一次情報（自分で収集）：アンケート・観察・実験・インタビュー
      \item 二次情報（既存のデータ）：書籍・論文・統計・Webサイト
    \end{itemize}
  \end{block}

  \vfill

  \begin{exampleblock}{補足}
    \begin{itemize}
      \item 複数の情報源を比較し，信頼性を確認する（評価5視点を活用）。出典を必ず記録する
    \end{itemize}
  \end{exampleblock}

  \vfill

\end{frame}

% --- Slide 4b: 演習（検索） ---

\begin{frame}{演習：インターネット検索}
  \begin{alertblock}{演習}
    \begin{enumerate}
      \item スマホで検索エンジンを開け
      \item 「情報セキュリティ」に関するAND検索を試せ
      \item OR検索，NOT検索もそれぞれ試せ
      \item 3パターンの検索式と結果の違いをノートにまとめよ
    \end{enumerate}
  \end{alertblock}
\end{frame}

% --- Slide 5: 発想・整理の手法 ---

\begin{frame}{発想・整理の手法（Theory 05）}
  \begin{block}{ブレーンストーミング}
    \begin{itemize}
      \item 複数人で自由にアイデアを出し合う手法
      \item 4つのルール：批判しない／自由に発言／質より量／他人の意見に便乗OK
    \end{itemize}
  \end{block}

  \vfill

  \begin{block}{KJ法}
    \begin{itemize}
      \item アイデアをカードに書き出し，グループ化して整理する手法
      \item 手順：カード作成 → グループ化 → 図解化 → 文章化
    \end{itemize}
  \end{block}

  \vfill

  \begin{block}{マインドマップ}
    \begin{itemize}
      \item 中心にテーマを置き，関連する語句を放射状に広げて整理する手法
      \item アイデアの全体像を視覚的に把握できる
    \end{itemize}
  \end{block}

  \vfill

\end{frame}

% --- Slide 5b: 演習（ブレスト・KJ法） ---

\begin{frame}{演習：ブレーンストーミング・KJ法}
  \begin{alertblock}{演習}
    \begin{enumerate}
      \item グループで学校生活に関するテーマを1つ決めよ（例：「通学時間の活用」「昼休みの過ごし方」）
      \item ブレーンストーミングで改善アイデアを10個以上出せ
      \item 付箋に1つずつ書き出せ
      \item KJ法の手順に従い，グループ分けして整理せよ
    \end{enumerate}
  \end{alertblock}
\end{frame}

% --- チェックリスト ---

\begin{frame}{まとめ・チェックリスト}
  \begin{block}{自己チェック（できたら $\checkmark$ をつけよう）}
    \begin{itemize}
      \item[$\square$] 問題解決の5ステップとPDCAサイクルを理解した
      \item[$\square$] インターネット検索のAND/OR/NOT検索を使えるようになった
      \item[$\square$] ブレーンストーミング・KJ法・マインドマップの発想・整理法を習得した
    \end{itemize}
  \end{block}

  \vfill

  {\small 質問があれば授業後またはTeamsで受け付けます。}
\end{frame}

\end{document}
